% !TEX root = ../main.tex
\appendix
\section{附录}

\section{GitHub 仓库地址}

本项目已开源至 GitHub，仓库地址为：

\begin{center}
\url{https://github.com/hxhdhcjmet/FoodShield.git}
\end{center}

仓库中包含了完整的系统源代码、攻击演示脚本、性能测试脚本和说明文档。读者可以克隆仓库到本地，按照 README 中的说明安装依赖并运行系统。

\section{项目目录结构}

FoodShield 项目采用模块化设计，主要目录结构如下：

\begin{verbatim}
FoodShield/
|-- project/
|   |-- crypto/                  # 国密算法、Token、PID、Merkle
|   |-- database/                # SQLite schema 与访问封装
|   |-- frontend/                # 用户端、骑手端、管理员端页面
|   |-- server/                  # Flask 主程序与审计逻辑
|   |-- tools/
|   |   |-- attack_demo.py       # 攻击/防御演示脚本
|   |   |-- performance_benchmark.py
|   |-- integration_demo.py      # 兼容入口，运行攻击/防御演示
|-- TODO.md
|-- requirements.txt
|-- pyproject.toml
|-- uv.lock
|-- foodshield.db                # 本地演示数据库
\end{verbatim}

\section{核心模块说明}

\subsection{国密算法模块（crypto/）}

\begin{itemize}[leftmargin=2em]
\item \texttt{sm\_utils.py}：SM3 哈希、HMAC-SM3 认证码、SM4-CBC 加解密的兼容封装。
\item \texttt{pid.py}：基于 HMAC-SM3 的匿名 PID 生成与验证。
\item \texttt{token\_utils.py}：订单 Token 生成与验证，绑定订单号、PID、时间戳和随机数。
\item \texttt{merkle.py}：消息哈希组织为 Merkle Tree，计算 Merkle Root 用于日志完整性验证。
\end{itemize}

\subsection{服务端模块（server/）}

\begin{itemize}[leftmargin=2em]
\item \texttt{app.py}：Flask 主程序，负责路由、WebSocket 房间管理、会话管理。
\item \texttt{logger.py}：Merkle 快照生成、完整性验证、异常定位。
\item \texttt{security\_audit.py}：数据库版安全审计辅助接口。
\end{itemize}

\subsection{前端页面（frontend/）}

\begin{itemize}[leftmargin=2em]
\item \texttt{user.html}：用户端页面，注册/登录、创建订单、Token 验证、聊天、订单历史、溯源申请。
\item \texttt{rider.html}：骑手端页面，骑手注册/登录、订单搜索、接单、匿名聊天、Root 上报、溯源申请。
\item \texttt{admin.html}：管理员端页面，登录、订单审计、消息哈希查看、Merkle 快照、完整性验证、三端一致性验证、溯源审批。
\end{itemize}

\section{攻击演示脚本说明}

攻击演示脚本位于 \texttt{project/tools/attack\_demo.py}，运行方式：

\begin{verbatim}
python -m project.tools.attack_demo
\end{verbatim}

脚本覆盖以下攻击场景：

\begin{enumerate}[label=\arabic*., leftmargin=2em]
\item 未登录创建订单被拒
\item 伪造 Token 验证失败
\item 其他用户越权验证订单被拒
\item 未登录骑手接单被拒
\item 重复接单被拒
\item 未登录管理员访问审计接口被拒
\item 管理员端消息接口只返回哈希，不返回 \texttt{content}
\item 直接篡改数据库消息密文后完整性验证失败
\item 管理员明文关键词扫描接口被隐私策略阻断
\end{enumerate}

\section{性能测试脚本说明}

性能测试脚本位于 \texttt{project/tools/performance\_benchmark.py}，运行方式：

\begin{verbatim}
python -m project.tools.performance_benchmark
\end{verbatim}

可调参数：

\begin{verbatim}
python -m project.tools.performance_benchmark --users 20 --orders-per-user 10 --messages-per-order 30
\end{verbatim}

脚本输出指标包括 Token 生成、Token 验证、SM4 加解密往返、消息写入、Merkle 快照和完整性验证的平均耗时、P95 和最大耗时。

\section{环境配置与启动说明}

\subsection{环境要求}

\begin{itemize}[leftmargin=2em]
\item Python 3.10+
\item 现代浏览器：Chrome、Edge 或 Firefox
\item 推荐使用虚拟环境或 uv
\end{itemize}

\subsection{安装与启动}

使用 pip：

\begin{verbatim}
python -m venv .venv
source .venv/bin/activate
pip install -r requirements.txt
python -m project.server.app
\end{verbatim}

Windows PowerShell 激活命令：

\begin{verbatim}
.\.venv\Scripts\Activate.ps1
\end{verbatim}

使用 uv：

\begin{verbatim}
uv sync
uv run python -m project.server.app
\end{verbatim}

启动后访问地址：

\begin{table}[H]
\centering
\caption{系统访问地址}
\label{tab:access-urls}
\begin{tabular}{ll}
\toprule
页面 & 地址 \\
\midrule
首页 & \texttt{http://127.0.0.1:5000/} \\
用户端 & \texttt{http://127.0.0.1:5000/user.html} \\
骑手端 & \texttt{http://127.0.0.1:5000/rider.html} \\
管理员端 & \texttt{http://127.0.0.1:5000/admin.html} \\
\bottomrule
\end{tabular}
\end{table}

\section{安全配置说明}

竞赛本地演示可使用默认配置。只要仍使用默认密钥或默认管理员密码，服务启动时会打印红色警告。公开网络演示或部署前应设置环境变量：

\begin{verbatim}
export FOODSHIELD_SECRET_KEY="replace-with-random-session-secret"
export FOODSHIELD_ADMIN_USERNAME="admin"
export FOODSHIELD_ADMIN_PASSWORD="replace-with-strong-password"
export FOODSHIELD_ORDER_KEY="replace-with-order-hmac-key"
export FOODSHIELD_MASTER_KEY="replace-with-pid-master-key"
export FOODSHIELD_ORDER_ALIAS_SECRET="replace-with-order-alias-key"
export FOODSHIELD_SM4_KEY="16-byte-sm4-key"
\end{verbatim}

\section{推荐演示流程}

\begin{enumerate}[label=\arabic*., leftmargin=2em]
\item 打开用户端，注册/登录，注意页面不展示 PID。
\item 创建订单，获得 \texttt{order\_id}、\texttt{timestamp} 和 \texttt{token}。
\item 验证订单 Token，进入匿名聊天房间。
\item 打开骑手端，注册或登录骑手账号。
\item 搜索待接订单，接单并进入同一订单房间。
\item 用户端和骑手端发送多条消息，观察 Merkle Root 同步。
\item 两端分别上报本端 Root。
\item 打开管理员端登录，查看订单列表和消息哈希。
\item 生成 Merkle 快照并执行完整性验证。
\item 执行三端 Root 一致性验证。
\item 用户端或骑手端提交溯源申请，管理员审批后按订单号溯源。
\item 运行攻击演示脚本，展示伪造 Token、越权、未登录骑手接单拦截、重复接单和篡改检测。
\item 运行性能测试脚本，收集报告中的性能数据。
\end{enumerate}
